\documentclass[12pt,a4paper]{exam}

% Paquetes necesarios
\usepackage[utf8]{inputenc}
\usepackage[T1]{fontenc}
\usepackage[spanish]{babel}
\usepackage{siunitx}
\usepackage{multicol}
\usepackage{geometry}

% Márgenes pequeños para meter 4 exámenes por hoja
\geometry{top=1cm,bottom=1cm,left=1cm,right=1cm}

\pagestyle{empty} % sin encabezados ni pies

\begin{document}

% ======== PRIMERA HOJA: Examen 1 y 2 ========
\begin{multicols}{2}

% ================= EXAMEN 1 ==================
\noindent
\makebox[\linewidth]{\textbf{Nombre y Apellido:}\enspace\hrulefill}
\vspace{1cm}
\begin{questions}
 \question Una máquina térmica con eficiencia $\eta_{\%}=87\%$ produce $W=4186\,J$ de trabajo. Encontrar: 
    \begin{parts}
        \part[2] El calor absorbido $Q_1$ del depósito caliente.
        \part[2] El calor desechado $Q_2$ al depósito frío.
    \end{parts} 
    
 \question En el funcionamiento de una máquina, se pierde el $30\%$ del calor entregado. Si la fuente caliente está a $T_1=\SI{400}{\celsius}$ y se obtiene un trabajo $W=200\,J$, calcular:
      \begin{parts}
          \part[1] El calor perdido $Q_2$. 
          \part[1] El rendimiento $\eta$.
          \part[1] El calor entregado $Q_1$.
          \part[2] Las temperaturas $T_1$ y $T_2$.
          \part[1] La nueva $T_1$ necesaria si se quiere duplicar el trabajo realizado. 
      \end{parts}
\end{questions}

\vfill
\columnbreak

% ================= EXAMEN 2 ==================
\noindent
\makebox[\linewidth]{\textbf{Nombre y Apellido:}\enspace\hrulefill}
\vspace{1cm}
\begin{questions}
 \question Una máquina térmica con eficiencia $\eta_{\%}=82\%$ produce $W=3500\,J$ de trabajo. Encontrar: 
    \begin{parts}
        \part[2] El calor absorbido $Q_1$ del depósito caliente.
        \part[2] El calor desechado $Q_2$ al depósito frío.
    \end{parts} 
    
 \question Una máquina pierde el $25\%$ del calor entregado. Si la fuente caliente está a $T_1=\SI{500}{\celsius}$ y el trabajo realizado es $W=250\,J$, calcular:
      \begin{parts}
          \part[1] El calor perdido $Q_2$. 
          \part[1] El rendimiento $\eta$.
          \part[1] El calor entregado $Q_1$.
          \part[2] Las temperaturas $T_1$ y $T_2$.
          \part[1] La nueva $T_1$ necesaria si se quiere duplicar el trabajo realizado. 
      \end{parts}
\end{questions}

\vfill
\end{multicols}


% ======== SEGUNDA HOJA: Examen 3 y 4 ========
\begin{multicols}{2}

% ================= EXAMEN 3 ==================
\noindent
\makebox[\linewidth]{\textbf{Nombre y Apellido:}\enspace\hrulefill}
\vspace{1cm}
\begin{questions}
 \question Una máquina térmica con eficiencia $\eta_{\%}=90\%$ produce $W=5000\,J$ de trabajo. Encontrar: 
    \begin{parts}
        \part[2] El calor absorbido $Q_1$ del depósito caliente.
        \part[2] El calor desechado $Q_2$ al depósito frío.
    \end{parts} 
    
 \question En una máquina se pierde el $35\%$ del calor entregado. Si la fuente caliente está a $T_1=\SI{450}{\celsius}$ y se realiza un trabajo de $W=300\,J$, calcular:
      \begin{parts}
          \part[1] El calor perdido $Q_2$. 
          \part[1] El rendimiento $\eta$.
          \part[1] El calor entregado $Q_1$.
          \part[2] Las temperaturas $T_1$ y $T_2$.
          \part[1] La nueva $T_1$ necesaria si se quiere duplicar el trabajo realizado. 
      \end{parts}
\end{questions}

\vfill
\columnbreak

% ================= EXAMEN 4 ==================
\noindent
\makebox[\linewidth]{\textbf{Nombre y Apellido:}\enspace\hrulefill}
\vspace{1cm}
\begin{questions}
 \question Una máquina térmica con eficiencia $\eta_{\%}=85\%$ produce $W=4500\,J$ de trabajo. Encontrar: 
    \begin{parts}
        \part[2] El calor absorbido $Q_1$ del depósito caliente.
        \part[2] El calor desechado $Q_2$ al depósito frío.
    \end{parts} 
    
 \question En el funcionamiento de una máquina, se pierde el $20\%$ del calor entregado. Si la fuente caliente está a $T_1=\SI{480}{\celsius}$ y el trabajo realizado es $W=220\,J$, calcular:
      \begin{parts}
          \part[1] El calor perdido $Q_2$. 
          \part[1] El rendimiento $\eta$.
          \part[1] El calor entregado $Q_1$.
          \part[2] Las temperaturas $T_1$ y $T_2$.
          \part[1] La nueva $T_1$ necesaria si se quiere duplicar el trabajo realizado. 
      \end{parts}
\end{questions}

\vfill
\end{multicols}

\end{document}
